\input{preamble.tex}
\title{Removing universal types in the HNN construction}
\author{Human Compression project}
\date{Private candidate, 11 October 2026}
\begin{document}\maketitle
The literal flag-transitivity assertion in Piedade--Tranchida,
arXiv:2505.24662v1, Theorem 3.35 is false on arbitrary nonfirm inputs,
as the separate reviewed $S_4$ example shows. This note recovers its
residually connected branch, including nonfirm inputs, by reducing to
the reviewed indexed HNN theorem. It does not assume that residual
connectedness means firmness.

Let $\alpha=(A,(A_i)_{i\in I})$ be a finite-rank flag-transitive,
residually connected coset geometry. Write $A_J=\bigcap_{i\in J}A_i$,
$H=A_I$, and $R_i=A_{I\setminus\{i\}}$.
The Borel-preserving residual generation formula is
\begin{equation}\label{eq:gen}
 A_J=\langle H,R_i:i\notin J\rangle.
\end{equation}
For zero or one remaining type it is the definition, and induction
on the complement size proves the general case from generation by
all next parabolics in each connected residue of rank at least two.
Thus no empty generated subgroup is substituted for $H$.

\begin{lemma}\label{lem:active}
Put $I_0=\{i:A_i=A\}$ and $I_+=I\setminus I_0$.
Then $A_i=A$ if and only if $R_i=H$. Deleting the universal types
$I_0$ leaves a firm flag-transitive, residually connected geometry
on $I_+$ with the same unlabelled family of standard parabolics and
the same Borel $H$. The original incidence geometry is its join
with one singleton for each type in $I_0$.
\end{lemma}
\begin{proof}
If $A_i=A$, then $H=A_i\cap R_i=R_i$. Conversely, if $R_i=H$,
formula (\ref{eq:gen}) for $J=\varnothing$ and $J=\{i\}$ gives
$A=A_i$. A type with maximal parabolic $A$ has exactly one element,
incident with every other element, proving the join description.
Dropping these factors leaves every standard intersection unchanged.
Flags still have common representatives. Every residue on at least
two active types is a residue of the original geometry after fixing
the deleted universal elements, so it remains connected. Each active
panel has index $[R_i:H]>1$, proving firmness. The empty active
geometry is interpreted as a single empty chamber.
\end{proof}

For later comparison we also make the firm branch explicit without
assuming residual connectedness. For a firm flag-transitive input
and a source-admissible $\phi:A_J\to A_K$,
the reviewed Boolean-lattice argument gives the indexed partial
bijection required by the HNN coset Helly theorem. For its classes
$s$ and $F=J\setminus K$, that theorem's intersection formula gives
the Borel $H$, the stable minimal parabolic $\langle H,t\rangle$,
and the other minimal parabolics
\[
 M_s=A_{I\setminus(s\setminus F)}.
\]
The first has infinite index over $H$, since the stable letter has
infinite image in $\ZZ$ and normalizes $H$. Every class has an
element outside $F$: a point in $F$ is in the range but not the
domain of the partial bijection, and its predecessor is outside
$J$, hence outside $F$. Consequently $M_s$ contains some
$R_i>H$. All panels are firm. This proves the full firm-input branch
of source Theorem 3.35, including part (c), independently of the
residual-connectivity assertion below.

Suppose $C=A_J,D=A_K$, and $\phi:C\to D$ satisfies the source's
Definition 3.30: it induces an isomorphism of the unlabelled families
of standard parabolics lying in these subgroups. The source forms
classes of the arrows $i\to j$, with $i\notin J,j\notin K$, for
which $\phi(A_{J\cup\{i\}})=A_{K\cup\{j\}}$. It uses
$F=J\setminus K$, the maximal subgroups $\langle A_s,t\rangle$
for these connected classes $s$, and the stable-type subgroup $A_F$
in the HNN extension $\langle A,t\mid t^{-1}ct=\phi(c)\rangle$.

\begin{theorem}\label{thm:rc}
With these literal choices, every flag-transitive, residually connected
input gives a flag-transitive, residually connected output. Thus the
residually connected branch of Theorem 3.35 remains valid without
firmness. Active output panels have the indices in the reviewed
indexed HNN theorem; each additional universal class has one element
in its panel.
\end{theorem}
\begin{proof}
Apply Lemma~\ref{lem:active}. On the active types the firm geometry
has the reverse Boolean parabolic lattice proved in the reviewed
indexed-admissibility note. Thus the isomorphism $\phi$ gives a
bijection
\[
 f:I_+\setminus J\longrightarrow I_+\setminus K
\]
with
$\phi(A_{J\cup X})=A_{K\cup f(X)}$ for every
$X\subseteq I_+\setminus J$. Here intersections with $J,K$
can ignore their inactive entries. This follows because the coatoms
of the Boolean interval inside $A_J$ are exactly the subgroups
$A_{J\cup\{i\}}$ for its missing active types; an isomorphism
permutes those coatoms and commutes with all their intersections.
The minimal subgroup $H$ is also preserved.

For an active $i\notin J$, the subgroup $A_{J\cup\{i\}}$ is
proper in $A_J$, whereas for an inactive $i$ it equals $A_J$.
The same distinction holds on the $K$ side. Hence a source arrow
never connects an active type to an inactive one. On active types
the connected classes are exactly those of the partial bijection
$f$, with isolated active types included. The stable parabolic is
$A_F=A_{F\cap I_+}$. Thus all active-class parabolics and the
stable parabolic coincide exactly with the reviewed indexed HNN
construction on the firm active geometry. Its flag transitivity and
residual connectedness follow from that theorem.

An inactive class $s$ consists only of types with $A_i=A$, so
$A_s=A$ and $\langle A_s,t\rangle$ is the entire HNN group.
Its coset type is a singleton universal element. Adding any number
of these types preserves flag transitivity and residual connectedness:
a residue with at least two active types uses the active conclusion,
a residue with one active type and at least one universal type has
a connected star, and a residue consisting of at least two universal
types is a complete graph. The same observation covers the case
$I_+=\varnothing$. In that case $C=D=A$ and the stable-type cosets
of $A$ in $A\rtimes\ZZ$ are the infinite rank-one part; all other
types are universal. This proves the assertion and panel description.
\end{proof}
The partial bijection and class compatibility, rather than a false
intersection formula for arbitrary special subgroups, are what justify
this conclusion. The reviewed nonfirm $S_4$ example is not residually
connected, so it does not contradict the theorem. This additional
positive reduction is a private candidate awaiting distinct review;
no formalization or source admission is inferred from it.
\end{document}
